\documentclass[14pt, a4paper]{article}
\usepackage{minitoc}
\usepackage[left=3.00cm, right=2.5cm, top=2.00cm, bottom=2.00cm]{geometry}
\usepackage{amsmath}
\usepackage{amssymb}
\usepackage{amsthm}
\usepackage{mathtools}
\usepackage{graphicx}
%\usepackage{algpseudocode}
%\usepackage{algorithm}
\usepackage[ruled,vlined,linesnumbered]{algorithm2e}
\usepackage{blindtext}
\usepackage{setspace}
\usepackage[utf8]{inputenc}
\usepackage[utf8]{vietnam}
\usepackage[center]{caption}
\usepackage[shortlabels]{enumitem}
\usepackage{fancyhdr} % header, footer
\usepackage{hyperref} % loại bỏ border với mục lục và công thức
\usepackage[nonumberlist, nopostdot, nogroupskip]{glossaries}
\usepackage{glossary-superragged}
\usepackage{tikz,tkz-tab}
\usepackage{pythonhighlight}
\setglossarystyle{superraggedheaderborder}
\pagestyle{fancy}
%\usepackage[style=numeric,sortcites]{biblatex}
%\addbibresource{ref.bib}
%\usepackage[numbers]{natbib}
\usepackage{indentfirst}
\usepackage[natbib,backend=biber,style=ieee, sorting=ynt]{biblatex}

\usepackage{caption}
\usepackage{subcaption}

\bibliography{ref.bib}

\graphicspath{{./figures/}}

\fancyhf{}
%\rhead{\textbf{Môn học: Các phương pháp thống kê hiện đại trong nghiên cứu Xã hội học}}
\lhead{\textbf{GVHD: TS. Trịnh Quốc Anh}}
\rfoot{\thepage}
\lfoot{\textbf{Học viên thực hiện: Nguyễn Chí Thanh - 21007925}}
\renewcommand{\headrulewidth}{0.4pt}
\renewcommand{\footrulewidth}{0.4pt}
%
%\numberwithin{equation}{section}
%\numberwithin{algorithm}{section}
%\numberwithin{figure}{section}
%
%\setlength{\parindent}{0.5cm}
%
%\setcounter{secnumdepth}{3} % Cho phép subsubsection trong report
%\setcounter{tocdepth}{3} % Chèn subsubsection vào bảng mục lục

%\newtheorem{dl}{Định lý}
%\newtheorem{md}{Mệnh đề}
%\newtheorem{bd}{Bổ đề}
%\newtheorem{dn}{Định nghĩa}
%\newtheorem{hq}{Hệ quả}

%\newtheorem{baitap}{Bài tập}
%\newtheorem*{loigiai}{Lời giải}

%\numberwithin{dl}{section}
%\numberwithin{md}{section}
%\numberwithin{bd}{section}
%\numberwithin{dn}{section}
%\numberwithin{hq}{section}

\setlength{\parindent}{0cm}

\newtheorem{dl}{Định lý}
\newtheoremstyle{sltheorem}
{}                % Space above
{}                % Space below
{\normalfont}        % Theorem body font % (default is "\upshape")
{}                % Indent amount
{\bfseries}       % Theorem head font % (default is \mdseries)
{.}               % Punctuation after theorem head % default: no punctuation
{ }               % Space after theorem head
{}                % Theorem head spec
\theoremstyle{sltheorem}
\newtheorem{baitap}{Bài tập}
\newtheoremstyle{soltheorem}
{}                % Space above
{}                % Space below
{\normalfont}        % Theorem body font % (default is "\upshape")
{}                % Indent amount
{\bfseries}       % Theorem head font % (default is \mdseries)
{.}               % Punctuation after theorem head % default: no punctuation
{\newline}               % Space after theorem head
{}                % Theorem head spec
\theoremstyle{soltheorem}
\newtheorem*{loigiai}{Lời giải}

\onehalfspacing

\begin{document}
\begin{titlepage}

    \newcommand{\HRule}{\rule{\linewidth}{0.5mm}} % Defines a new command for the horizontal lines, change thickness here

    \center % Center everything on the page

    %----------------------------------------------------------------------------------------
    %	HEADING SECTIONS
    %----------------------------------------------------------------------------------------
    \textsc{\LARGE Đại học Quốc Gia Hà Nội}\\[0.5cm]
    \textsc{\LARGE Trường đại học Khoa học tự nhiên}\\[0.5cm] % Name of your university/college
    \textsc{\LARGE Khoa Toán - Cơ - Tin học}\\[0.5cm]

    \includegraphics[scale=0.2]{HUS-logo.jpg}\\[0.5cm]

    \textsc{\Large Chuyên ngành: Khoa học dữ liệu}\\[0.5cm] % Major heading such as course name


    %----------------------------------------------------------------------------------------
    %	TITLE SECTION
    %----------------------------------------------------------------------------------------

    \HRule \\[0.4cm]
    { \huge \bfseries Bài tập môn học}\\[0.4cm] % Title of your document
    \HRule \\[1.5cm]

    \textsc{\Large Môn học: Các phương pháp thống kê hiện đại \\ trong nghiên cứu Xã hội học}\\[1cm] % Minor heading such as course title


    \textsc{\Large Bài tập số 3}\\[1cm]


    %----------------------------------------------------------------------------------------
    %	AUTHOR SECTION
    %----------------------------------------------------------------------------------------
    \begin{minipage}{0.4\textwidth}
        \begin{flushleft} \large
        \emph{Giảng viên hướng dẫn:} \\
        TS. Trịnh Quốc Anh % Supervisor's Name
        \end{flushleft}
    \end{minipage}\\[0.5cm]

    \begin{minipage}{0.4\textwidth}
    \begin{flushleft} \large
    \emph{Học viên thực hiện:}\\
    Nguyễn Chí Thanh \\
    MSHV: 21007925 \\ % Your name
    Lớp: Khoa học dữ liệu - K4
    \end{flushleft}
    \end{minipage}


    % If you don't want a supervisor, uncomment the two lines below and remove the section above
    %\Large \emph{Author:}\\
    %John \textsc{Smith}\\[3cm] % Your name

    %----------------------------------------------------------------------------------------
    %	DATE SECTION
    %----------------------------------------------------------------------------------------

    % I don't want day because it is English
    % {\large \today}\\[2cm] % Date, change the \today to a set date if you want to be precise

    %----------------------------------------------------------------------------------------
    %	LOGO SECTION
    %----------------------------------------------------------------------------------------

    %\includegraphics{logo/rsz_3logo-khtn.png}\\[1cm] % Include a department/university logo - this will require the graphicx package

    %----------------------------------------------------------------------------------------

    \vfill % Fill the rest of the page with whitespace

\end{titlepage}

\nocite{*}

\newpage

\begin{baitap}
    Áp dụng các phương pháp Ridge, LASSO, và PCA cho dữ liệu Sonar.

    Hãy so sánh các kết quả thu được và biện luận để đưa ra mô hình phù hợp nhất.
\end{baitap}

\begin{loigiai}
    Hàm mất mát của mô hình logistic regression:
    \begin{equation*}
        \mathcal{L} = -\dfrac{1}{N}\sum_{i=1}^N (y_i \log \hat{y}_i + (1-y_i)\log(1-\hat{y}_i)) 
    \end{equation*}

    với $\hat{y}=p(y=1\vert \bold{x})=f(\bold{x}) = \sigma(w_0 + w_1 x_1 + w_2 x_2 + \dots w_n x_n) = \sigma(\bold{w}^T \bold{x})$ là kết quả đầu ra của mô hình logistic tương ứng với đầu vào $\bold{x}$.

    Hàm mất mát của phương pháp Ridge:

    \begin{equation*}
        \mathcal{L}_{\mathrm{Ridge}} = -\dfrac{1}{N}\sum_{i=1}^N (y_i \log \hat{y}_i + (1-y_i)\log(1-\hat{y}_i)) + \lambda \lVert \bold{w} \rVert_2^2
    \end{equation*}

    Hàm mất mát của phương pháp LASSO:

    \begin{equation*}
        \mathcal{L}_{\mathrm{Ridge}} = -\dfrac{1}{N}\sum_{i=1}^N (y_i \log \hat{y}_i + (1-y_i)\log(1-\hat{y}_i)) + \lambda \lVert \bold{w} \rVert_1
    \end{equation*}

    Ta sử dụng 5 mô hình đã được sử dụng trong bài tập trước:

    \begin{itemize}
        \item Mô hình gồm tất cả các biến
        \item Mô hình forward được chọn ở bài tập trước
        \item Mô hình backward được chọn ở bài tập trước
        \item Mô hình forward-backward được chọn ở bài tập trước
        \item Mô hình các biến được trích xuất từ thuật toán PCA giữ lại 95 \% thông tin
    \end{itemize}

    Để chọn ra mô hình tốt nhất, mỗi mô hình đi kèm với cách chỉnh định sẽ được tối ưu hóa các siêu tham số gồm:

    \begin{itemize}
        \item Phương pháp tiền xử lý: min-max scaling, standardize.
        \item Tốc độ học được biến đổi trong quá trình huấn luyện theo các cách: fixed, cyclic, inverse\_decay, exponential\_decay, factor\_decay, squareroot, cosine.
        \item Các tốc độ học khác nhau: 1e-4, 1e-3, 1e-2, 1e-1, 1, 2, 5, 10
        \item Các thuật toán tối ưu hóa, 12 thuật toán tối ưu hóa được sử dụng gồm:
        \begin{itemize}
            \item Gradient Descent
            \item Adam
            \item Avagrad
            \item RAdam
            \item Momentum
            \item Adagrad
            \item RMSProp
            \item Adadelta
            \item Adamax
            \item Nadam
            \item AMSGrad
            \item AdaBelief
        \end{itemize}

        Sơ đồ mã giả của các thuật toán tối ưu hóa được trình bày ở phụ lục \ref{Optimization-Algorithms}
    \end{itemize}

    Các siêu tham số trên sẽ được huấn luyện sử dụng code python và numpy không sử dụng thư viện khác.
    Việc sử dụng code python và numpy để có thể thử với các thuật toán tối ưu hóa khác mà thư viện sklearn không có.

    Ngoài ra ta vẫn huấn luyện một tập các mô hình được huấn luyện sử dụng thư viện sklearn.
    Siêu tham số của các mô hình được huấn luyện sử dụng sklearn bao gồm:

    \begin{itemize}
        \item Phương pháp tiền xử lý: min-max scaling, standardize.
        \item Các hệ số $C$ là nghịch đảo của $\lambda$
        \item Các thuật toán tối ưu hóa, 12 thuật toán tối ưu hóa được sử dụng gồm:
    \end{itemize}

    Mỗi một cặp thông tin gồm phương pháp chỉnh định ("l1", "l2", không chỉnh định) và mô hình sẽ được chọn ra bộ siêu tham số tốt nhất từ các siêu tham số được đề cập ở trên bằng độ đo AUC-ROC trên tập test.
    Ta sẽ cố định tập test ngay từ đầu.
    Như vậy ta sẽ có 5 (mô hình) $\times$ 3 ("l1"-lasso, "l2"-ridge, không chỉnh định) = 15 mô hình ứng viên được chọn từ quá trình tối ưu hóa siêu tham số.

    Để chọn ra mô hình tốt nhất từ nhiều yếu tố như $AUC-ROC, AIC, BIC$ cũng như từ ma trận nhầm lẫn ứng với ngưỡng xác suất $p^*$ được chọn.



\end{loigiai}

\newpage
\printbibliography[title={TÀI LIỆU THAM KHẢO}]

\newpage
\appendix

\section{Các thuật toán tối ưu hóa} \label{Optimization-Algorithms}


\begin{algorithm}[h!]
    \DontPrintSemicolon
    \KwIn{độ dài bước $\lbrace \alpha_t \rbrace_{t=1}^{T}$, $w_0$ khởi tạo, hàm mục tiêu $\ell(w)$}
    \KwOut{$w$ đã được học}

    \For{$t \gets 1$ \KwSty{to} $T$} {
        $g_t \gets \nabla \ell(w_{t-1})$\;
        $w_t \gets w_{t-1} - \alpha_t g_t$
    }
    \Return{$w_T$}\;
    \caption{Thuật toán Gradient Descent}
\end{algorithm}


\begin{algorithm}[h!]
    \DontPrintSemicolon
    \KwIn{độ dài bước $\lbrace \alpha_t \rbrace_{t=1}^{T}$, hệ số $\beta_1$, $w_0$ khởi tạo, hàm mục tiêu $\ell(w)$}
    \KwOut{$w$ đã được học}
    $m_0 \gets 0$\;
    \For{$t \gets 1$ \KwSty{to} $T$} {
        $g_t \gets \nabla \ell(w_{t-1})$\;
        $m_t \gets \beta_1 m_{t-1} + (1-\beta_1) g_t$\;
        $w_t \gets w_{t-1} - \alpha_t m_t$\;
    }
    \Return{$w_T$}\;
    \caption{Thuật toán Momentum}
\end{algorithm}


\begin{algorithm}[h!]
    \KwIn{độ dài bước $\lbrace \alpha_t \rbrace_{t=1}^{T}$, $w_0$ khởi tạo, hàm mục tiêu $\ell(w)$}
    \KwOut{$w$ đã được học}

    $v_0 \gets 0$\;
    \For{$t \gets 1$ \KwSty{to} $T$} {
        $g_t \gets \nabla \ell(w_{t-1})$\;
        $v_t \gets v_{t-1} + g_t^2$\;
        $w_t \gets w_{t-1} - \alpha_t \dfrac{g_t}{\sqrt{v_t} + \epsilon}$\;
    }

    \Return{$w_T$}\;
    \caption{Thuật toán Adagrad}
\end{algorithm}

\begin{algorithm}[h!]
    \DontPrintSemicolon
    \KwIn{độ dài bước $\lbrace \alpha_t \rbrace_{t=1}^{T}$, hệ số $\beta$, $w_0$ khởi tạo, hàm mục tiêu $\ell(w)$}
    \KwOut{$w$ đã được học}

    $v_0 \gets 0$\;
    $d_0 \gets 0$\;

    \For{$t \gets 1$ \KwSty{to} $T$} {
        $g_t \gets \nabla \ell (w_{t-1})$\;
        $v_t \gets \beta v_{t-1} + (1-\beta) g_t^2$\;
        $\Delta w \gets -\alpha_t \dfrac{\sqrt{d_{t-1} + \epsilon}g_t}{\sqrt{v_t + \epsilon}}$\;
        $w_t \gets w_{t-1} + \delta w$\;
        $d_t \gets \beta d_{t-1} + (1-\beta) \delta w^2$\;
    }

    \Return{$w_T$}\;
    \caption{Thuật toán Adadelta}
\end{algorithm}


\begin{algorithm}[h!]
    \DontPrintSemicolon
    \KwIn{độ dài bước $\lbrace \alpha_t \rbrace_{t=1}^{T}$, hệ số $\beta$, $w_0$ khởi tạo, hàm mục tiêu $\ell(w)$}
    \KwOut{$w$ đã được học}

    $v_0 \gets 0$\;
    \For{$t \gets 1$ \KwSty{to} $T$} {
        $g_t \gets \nabla \ell (w_{t-1})$\;
        $v_t \gets \beta v_{t-1} + (1-\beta) g_t^2$\;
        $w_t \gets w_{t-1} - \alpha_t \dfrac{g_t}{\sqrt{v_t} + \epsilon}$\;
    }

    \Return{$w_T$}\;
    \caption{Thuật toán RMSProp}
\end{algorithm}


\begin{algorithm}[h!]
    \DontPrintSemicolon
    \KwIn{độ dài bước $\lbrace \alpha_t \rbrace_{t=1}^{T}$, các hệ số $\beta_1, \beta_2$, $w_0$ khởi tạo, hàm mục tiêu $\ell(w)$}
    \KwOut{$w$ đã được học}

    $m_0 \gets 0$\;
    $v_0 \gets 0$\;

    \For{$t \gets 1$ \KwSty{to} $T$} {
        $g_t \gets \nabla \ell(w_{t-1})$\;
        $m_t \gets \beta_1 m_{t-1} + (1 - \beta_1)g_t$\;
        $v_t \gets \beta_2 v_{t-1} + (1 - \beta_2)g_t^2$\;
        $\hat{m}_t \gets \dfrac{m_t}{1 - \beta_1^t}$\;
        $\hat{v}_t \gets \dfrac{v_t}{1 - \beta_2^t}$\;
        $w_t \gets w_{t-1} - \alpha_t \dfrac{\hat{m}_t}{\sqrt{\hat{v}_t} + \epsilon}$\;
    }

    \Return{$w_T$}\;
    \caption{Thuật toán Adam}
\end{algorithm}


\begin{algorithm}[h!]
    \DontPrintSemicolon
    \KwIn{độ dài bước $\lbrace \alpha_t \rbrace_{t=1}^{T}$, các hệ số $\beta_1, \beta_2$, $w_0$ khởi tạo, hàm mục tiêu $\ell(w)$}
    \KwOut{$w$ đã được học}

    $m_0 \gets 0$\;
    $u_0 \gets 0$\;

    \For{$t \gets 1$ \KwSty{to} $T$} {
        $g_t \gets \nabla \ell(w_{t-1})$\;
        $m_t \gets \beta_1 m_{t-1} + (1 - \beta_1)g_t$\;
        $u_t \gets \max (\beta u_{t-1}, \lVert g_t \rVert_{\infty})$\;
        $w_t \gets w_{t-1} - \alpha_t \dfrac{m_t}{(1-\beta_1^t)u_t + \epsilon}$\;
    }

    \Return{$w_T$}\;
    \caption{Thuật toán Adamax}
\end{algorithm}


\begin{algorithm}[h!]
    \DontPrintSemicolon
    \KwIn{độ dài bước $\lbrace \alpha_t \rbrace_{t=1}^{T}$, các hệ số $\beta_1, \beta_2$, $w_0$ khởi tạo, hàm mục tiêu $\ell(w)$}
    \KwOut{$w$ đã được học}

    $m_0 \gets 0$\;
    $v_0 \gets 0$\;

    \For{$t \gets 1$ \KwSty{to} $T$} {
        $g_t \gets \nabla \ell(w_{t-1})$\;
        $m_t \gets \beta_1 m_{t-1} + (1 - \beta_1)g_t$\;
        $v_t \gets \beta_2 v_{t-1} + (1 - \beta_2)g_t^2$\;
        $\hat{v}_t \gets \max(\hat{v}_{t-1}, v_t)$\;
        $w_t \gets w_{t-1} - \alpha_t \dfrac{m_t}{\sqrt{\hat{v}_t} + \epsilon}$\;
    }

    \Return{$w_T$}\;
    \caption{Thuật toán AMSGrad}
\end{algorithm}


\begin{algorithm}[h!]
    \DontPrintSemicolon
    \KwIn{độ dài bước $\lbrace \alpha_t \rbrace_{t=1}^{T}$, các hệ số $\beta_1, \beta_2$, $w_0$ khởi tạo, hàm mục tiêu $\ell(w)$}
    \KwOut{$w$ đã được học}

    $m_0 \gets 0$\;
    $v_0 \gets 0$\;

    \For{$t \gets 1$ \KwSty{to} $T$} {
        $g_t \gets \nabla \ell(w_{t-1})$\;
        $m_t \gets \beta_1 m_{t-1} + (1 - \beta_1)g_t$\;
        $v_t \gets \beta_2 v_{t-1} + (1 - \beta_2)g_t^2$\;
        $\hat{m}_t \gets \dfrac{m_t}{1 - \beta_1^t}$\;
        $\hat{v}_t \gets \dfrac{v_t}{1 - \beta_2^t}$\;
        $w_t \gets w_{t-1} - \dfrac{\alpha_t}{\sqrt{\hat{v}_t} + \epsilon} \Bigg( \beta_1 \hat{m}_t + \dfrac{1 - \beta_1}{1 - \beta_1^t}g_t \Bigg)$\;
    }

    \Return{$w_T$}\;
    \caption{Thuật toán Nadam}
\end{algorithm}


\begin{algorithm}[h!]
    \DontPrintSemicolon
    \KwIn{độ dài bước $\lbrace \alpha_t \rbrace_{t=1}^{T}$, các hệ số $\beta_1, \beta_2$, $w_0$ khởi tạo, hàm mục tiêu $\ell(w)$}
    \KwOut{$w$ đã được học}

    $m_0 \gets 0$\;
    $s_0 \gets 0$\;

    \For{$t \gets 1$ \KwSty{to} $T$} {
        $g_t \gets \nabla \ell(w_{t-1})$\;
        $m_t \gets \beta_1 m_{t-1} + (1 - \beta_1)g_t$\;
        $s_t \gets \beta_2 s_{t-1} + (1-\beta_2)(g_t - m_t)^2 + \epsilon$\;
        $\hat{m}_t \gets \dfrac{m_t}{1 - \beta_1^t}$\;
        $\hat{s}_t \gets \dfrac{s_t}{1 - \beta_2^2}$\;
        $w_t \gets w_{t-1} - \alpha_t \dfrac{\hat{m}_t}{\sqrt{\hat{s}_t} + \epsilon}$
    }

    \Return{$w_T$}\;
    \caption{Thuật toán AdaBelief}
\end{algorithm}


\begin{algorithm}[h!]
    \DontPrintSemicolon
    \KwIn{độ dài bước $\lbrace \alpha_t \rbrace_{t=1}^{T}$, các hệ số $\beta_1, \beta_2$, $w_0$ khởi tạo, hàm mục tiêu $\ell(w)$}
    \KwOut{$w$ đã được học}
    $m_0 \gets 0$\;
    $v_0 \gets 0$\;
    $\rho_{\infty} \gets 2/(1-\beta_2)-1$\;
    \For{$t \gets 1$ \KwSty{to} $T$}{
        $g_t \gets \nabla \ell(w_{t-1})$\;
        $m_t \gets \beta_1 m_{t-1} + (1 - \beta_1)g_t$\;
        $v_t \gets \beta_2 v_{t-1} + (1 - \beta_2)g_t^2$\;
        $\widehat{m}_t \gets m_t / (1 - \beta_1^t)$\;
        $\rho_t \gets \rho_{\infty} - 2t\beta_2^t / (1 - \beta_2^t)$\;
        \If {$\rho_t < 4$} {
            $\widehat{v}_t \gets \sqrt{v_t / (1 - \beta_2^t)}$\;
            $r_t \gets \sqrt{\frac{(\rho_t - 4)(\rho_t - 2)\rho_{\infty}}{(\rho_{\infty} - 4)(\rho_{\infty} - 2)\rho_t}}$\;
            $w_t \gets w_{t-1} - \alpha_t r_t \widehat{m}_t / (\widehat{v}_t + \epsilon)$\;
        } \Else {
            $w_t \gets w_{t-1} - \alpha_t \widehat{m}_t$\;
        }
    }
    \Return{$w_T$}\;
    \caption{Thuật toán RAdam}
\end{algorithm}


\begin{algorithm}[h!]
    \DontPrintSemicolon
    \KwIn{độ dài bước $\lbrace \alpha_t \rbrace_{t=1}^{T}$, hệ số $\beta_1$, $w_0$ khởi tạo, hàm mục tiêu $\ell(w)$}
    \KwOut{$w$ đã được học}
    $m_0 \gets 0$\;
    $v_0 \gets 0$\;

    \For{$t \gets 1$ \KwSty{to} $T$}{
        $g_t \gets \nabla \ell(w_{t-1})$\;
        $m_t \gets \beta_1 m_{t-1} + (1 - \beta_1)g_t$\;
        $\eta_t \gets \dfrac{1}{\sqrt{v_{t-1}} + \epsilon}$\;
        $w_t \gets w_{t-1} - \alpha_t \dfrac{\eta_t}{\lVert \eta_t / \sqrt{d} \rVert_2 \odot m_t}$\;
        $v_t \gets \beta_2 v_{t-1} + (1-\beta_2)g_t^2$\;
    }
    \Return{$w_T$}\;
    \caption{Thuật toán Avagrad}
\end{algorithm}

\section{Các phương pháp điều chỉnh tốc độ học}


\begin{algorithm}[h!]
    \DontPrintSemicolon
    \KwIn{Tốc độ học ban đầu $\alpha$}
    \For{$t \gets 0$ \KwSty{to} $T$} {
        $t \gets t + 1$\;
        $\alpha_t \gets \alpha$\;
    }
    \caption{Điều chỉnh tốc độ học theo phương pháp cố định}
\end{algorithm}


\begin{algorithm}[h!]
    \DontPrintSemicolon
    \KwIn{Tốc độ học ban đầu $\alpha$}
    \For{$t \gets 0$ \KwSty{to} $T$} {
        $t \gets t + 1$\;
        $\alpha_t \gets \dfrac{\alpha}{t}$\;
    }
    \caption{Điều chỉnh tốc độ học theo phương pháp inverse decay}
\end{algorithm}


\begin{algorithm}[h!]
    \DontPrintSemicolon
    \KwIn{min\_lr, max\_lr, num\_increase, num\_decrease}
    \For{$t \gets 0$ \KwSty{to} $T$} {
        $t \gets t + 1$\;
        res\_step $\gets$ num\_steps mod (num\_increase + num\_decrease)\;
        \If{res\_step < num\_increase} {
            $\alpha_t \gets$  min\_lr + (max\_lr - min\_lr) $\times$ res\_step / num\_increase\;
        } \Else {
            $\alpha_t \gets$ max\_lr - (max\_lr - min\_lr) $\times$ (res\_step - num\_increase) / (num\_decrease)\;
        }
    }
    \caption{Điều chỉnh tốc độ học theo phương pháp cyclic}
\end{algorithm}


\begin{algorithm}[h!]
    \DontPrintSemicolon
    \KwIn{min\_lr, max\_lr, max\_steps}
    \For{$t \gets 0$ \KwSty{to} $T$} {
        $t \gets t + 1$\;
        \If{$t \geq$ max\_steps} {
            $\alpha_t \gets$ min\_lr\;
        } \Else {
            $\alpha_t \gets$ max\_lr - num\_steps $\times$ (init\_lr - min\_lr) / max\_steps\;
        }
    }
    \caption{Điều chỉnh tốc độ học theo phương pháp giảm tuyến tính}
\end{algorithm}

\begin{algorithm}[h!]
    \DontPrintSemicolon
    \KwIn{Tốc độ học ban đầu $\alpha$, tốc độ giảm $k$}
    \For{$t \gets 0$ \KwSty{to} $T$} {
        $t \gets t + 1$\;
        $\alpha_t \gets \alpha \exp(-k \times t)$\;
    }
    \caption{Điều chỉnh tốc độ học theo phương pháp giảm theo hàm mũ}
\end{algorithm}

\begin{algorithm}[h!]
    \DontPrintSemicolon
    \KwIn{Tốc độ học ban đầu $\alpha$, hệ số giảm $k < 1$, chu kỳ giảm $T$}
    \For{$t \gets 0$ \KwSty{to} $T$} {
        $t \gets t + 1$\;
        \If {$t \% T = 0$} {
            $\alpha_t \gets \alpha \times k$\;
        } \Else {
            $\alpha_t \gets \alpha_{t-1}$\;
        }
    }
    \caption{Điều chỉnh tốc độ học theo phương pháp giảm theo hàm mũ}
\end{algorithm}

\begin{algorithm}[h!]
    \DontPrintSemicolon
    \KwIn{Tốc độ học ban đầu $\alpha$}
    \For{$t \gets 0$ \KwSty{to} $T$} {
        $t \gets t + 1$\;
        $\alpha_t \gets \dfrac{\alpha}{\sqrt{t}}$\;
    }
    \caption{Điều chỉnh tốc độ học theo phương pháp giảm theo hàm căn}
\end{algorithm}


\begin{algorithm}[h!]
    \DontPrintSemicolon
    \KwIn{min\_lr, max\_lr, max\_steps}
    \For{$t \gets 0$ \KwSty{to} $T$} {
        $t \gets t + 1$\;
        $\alpha_t \gets$ min\_lr + $\dfrac{1}{2}$ (max\_lr - min\_lr) (1 + $\cos$((num\_steps $\times \pi$)/max\_steps))\;
    }
    \caption{Điều chỉnh tốc độ học theo phương pháp cosine}
\end{algorithm}


\end{document}